% ===========================================================================
\section{Limitations}
\label{sec:limitations}
% ===========================================================================

The estimates are provisional and desk-based. Several limitations follow directly
from the method and should condition interpretation of the results.

\subsection*{Emission factors and tier status}
\begin{enumerate}[leftmargin=1.4em,itemsep=2pt,topsep=2pt]

\item \textbf{Tier~1/2, not Tier~3.} No emission factor is backed by Pakistani
field measurement; the factors are IPCC defaults elevated with country-specific
activity data. The widest uncertainty bands --- rice and soil nitrous oxide ---
reflect this, and it is the principal limitation on inventory confidence.

\item \textbf{Emission-factor context.} With the exception of the empirical
soil-nitrogen anchor, factors originate from IPCC regional defaults; combustion,
manure, and soil processes in Pakistani conditions may differ systematically from
the values underlying those defaults, in either direction.

\item \textbf{The soil direct factor is a deliberate choice.} We use the 2019
dry-climate $\EFone{}=0.005$, which halves managed soils relative to the 2006
aggregate $0.010$ that FAOSTAT retains (Section~\ref{sec:discussion}). We report
the aggregate value as a sensitivity ($\sim$\SI{40}{}~Mt) rather than tune to the
benchmark; the choice is defensible but consequential.

\item \textbf{Rice is modelled as continuously flooded} ($SF_w=1$). The
intermittent and rainfed water-regime scaling factors exist in the factor set but
are not applied, so alternate-wetting-and-drying mitigation is not represented ---
a real simplification on a sizeable, policy-relevant category with a wide
uncertainty band.

\end{enumerate}

\subsection*{Livestock and manure management}
\begin{enumerate}[leftmargin=1.4em,itemsep=2pt,topsep=2pt,resume]

\item \textbf{Manure-management shares are uncertain.} The management-system
split is Tier-2/FAOSTAT-consistent; an alternative South-Asian split would move
manure direct nitrous oxide materially, and the split is not measured for
Pakistan.

\item \textbf{The dung-cake transfer has an observed destination for part of its mass.}
Dung-cake combustion ($\sim$28\% of manure) is removed from this sector's
manure-management balance. The national reference books as household fuel what the
household survey observes, under a fifth of the diverted mass; the remainder is carried
as an unresolved row outside the total with no emission applied, and its candidate fates
are priced rather than chosen. Under national configuration~L, the reference of earlier
releases and now the alternative, all of the diverted mass is RECEIVED by the residential
sector as household fuel, and there the cross-module closure is enforced with an exact
equality, not merely declared (an earlier edition of this item still called the
receiving account out of scope).

\item \textbf{The crop-residue fuel transfer is closed the same way, with a larger
open remainder.} Of \SI{86.14}{}~Mt of dry residue, \SI{9.27}{}~Mt (\SI{122192}{}~TJ) is
diverted to household fuel and received by the residential sector's fuel ledger under the
same exact-equality discipline as dung cake; \SI{20.40}{}~Mt is returned to soil and
\SI{8.37}{}~Mt burned in the field. The \SI{48.10}{}~Mt the ledger cannot place
(\SI{55.8}{\percent} of dry residue---fodder, bedding, industrial use and decay) is
carried as an explicit unresolved row rather than assigned a fate.

\item \textbf{Livestock ammonia is \boldmath$\sim$3$\times$ independent
inventories, and is provisional.} Benchmarking (Section~\ref{sec:ap-benchmark})
shows our synthetic-fertiliser \NHthree{} at 1.20$\times$ and the livestock streams at
2.5$\times$ an independent South Asian inventory \citep{sanh2026}, both compared at the
benchmark's own year. The shared nitrogen ledger
(Section~\ref{sec:results-nledger}) makes the loss formulation a secondary rather than
a primary suspect, but does not eliminate it: the IPCC and EMEP livestock totals agree
to within 1.23$\times$ on a common basis, yet field application and grazing disagree by
factors of 0.66 and 3.5 in opposite directions and partly offset. What it does \emph{not}
do is establish which side is right, because both representations act on the same
excreted-nitrogen pool --- raised by our Tier-2 animal masses, which also leave
manure management at 1.38$\times$ FAOSTAT --- and the benchmark's own livestock
parameters are not fully reproducible (its implicit-emission-factor table omits
buffalo, the dominant species here). We therefore treat the livestock \NHthree{}
estimate and the manure-derived \NtwoO{} that shares its nitrogen as provisional.
Matching the benchmark's \emph{year} removed roughly a third of the apparent gap,
and back-calculating its unpublished buffalo factor
(Section~\ref{sec:ap-benchmark}) points most of the remainder towards the
manure-management split rather than the animal masses or the loss chain. That split
is not measured for Pakistan, and the two external parameterisations that bear on it --- GLEAM 3.0
and this benchmark --- bracket our assumption from opposite sides and disagree
with each other by a factor of three. A national manure-management survey is therefore
the single highest-value data collection identified by this work, and it would move
both the \NHthree{} and the manure \COe{} estimates.

\end{enumerate}

\subsection*{Historical and spatial assumptions}
\begin{enumerate}[leftmargin=1.4em,itemsep=2pt,topsep=2pt,resume]

\item \textbf{Historical source and time-basis transition.} The 1990--2023 series
uses FAOSTAT activity data on a calendar-year basis while 2024 is anchored to the
PBS census on the FY2023--24 fiscal basis; no statistical harmonisation is applied
across this join---neither for the source change nor for the six-month basis
offset---so part of the 2023--2024 change may
reflect either transition rather than true annual activity growth (livestock
species are mapped to a common six-species taxonomy, but levels can differ). Crop
coverage also differs: the historical file carries four crops where the 2024 anchor
carries five, so residue-driven quantities --- the burning greenhouse gases and the
burning air pollutants --- are a few per cent lower before the join than after it.

\item \textbf{Gridded and district maps are allocated surfaces, not
observations.} Within each province the pattern follows a proxy; the
reconciliation invariant guarantees the totals but not the sub-provincial detail.

\item \textbf{District spatial shares are held constant across years.} The
district\,$\times$\,year panel (Section~\ref{sec:results-districtyear}) rescales a
single 2024 spatial pattern to each year's national total, because no multi-year
district activity survey exists. It distributes national growth across space; it
does not detect district-specific shifts in activity.

\end{enumerate}

\subsection*{Validation and uncertainty}
\begin{enumerate}[leftmargin=1.4em,itemsep=2pt,topsep=2pt,resume]

\item \textbf{District validation is uneven.} Only rice and (via wheat area)
managed soils have an independent district record; the livestock, burning and
urea-application categories have none, and their district rankings are therefore unvalidated
against official statistics (Section~\ref{sec:validation-district}).

\item \textbf{Uncertainty treats most components as independent.} Perfect correlation
is modelled among terms that share an emission factor within managed soils --- the
indirect volatilisation and leaching streams, which are separate rows only because
they follow different spatial proxies. All other correlations (shared activity data,
common factor families across categories) are omitted \emph{here}, because this stream
samples component totals rather than parameters, and could widen or narrow the
national interval. The national engine samples a parameter graph instead and carries two
of these couplings explicitly; Supplementary Note~S3 gives their magnitude.

\item \textbf{The reported ammonia interval is parametric, not total.} The
Monte-Carlo band holds the manure-management configuration fixed and varies factors
and activity within it. But that configuration is itself the dominant uncertainty:
across defensible alternatives national \NHthree{} spans \SI{1.67}{}--\SI{2.67}{}~Mt
(Section~\ref{sec:ap-structural}), a wider range than the $\pm$25\% parametric band.
The two are reported separately and should not be combined into a single interval,
because the structural term is a choice between discrete configurations rather than a
distribution.

\item \textbf{Air-pollutant totals are bottom-up, not control totals.} The
air-pollutant layer now carries the same gridded, district, time-series,
Monte-Carlo and key-source treatment as the greenhouse gases, and its spatial
products reconcile to its national totals by construction. But those national
totals are our own bottom-up estimates rather than an audited IPCC control total,
so what the allocation conserves is weaker in provenance than for the GHGs. The
crop-burning factors are IPCC/EMEP/biomass-burning defaults, soil \NOx{} is a small
high-uncertainty term, and none of the layer is field-validated here.

\end{enumerate}

\subsection*{Scope and omitted sources}
\begin{enumerate}[leftmargin=1.4em,itemsep=2pt,topsep=2pt,resume]

\item \textbf{Not-estimated agriculture gaps.} Soil-organic-nitrogen mineralisation
($F_{SOM}$) and liming are marked NE in the completeness table
(\ref{app:completeness}); they require field or national activity data. (Land-sector
gaps --- SOC change, harvested wood products --- are covered in Supplementary
Note~S1.)

\item \textbf{Land sector excluded.} The exploratory Tier-1 land sector is the
least settled part of this work (single-year, contested mangrove sign) and is
reported only in Supplementary Note~S1, excluded from all agricultural results.

\item \textbf{The fire layer is an observation, and observations can be absent.} The
burning surface now rides the reporting period's own detections (the offset window
that stood here---October 2024--May 2025, one season late---was replaced by the
period-matched extraction). The honest cost is sparsity: a district shape carrying
residue-burning panel mass in which VIIRS recorded no fire falls back to a uniform
spread within that shape, declared per event in the grid fallback ledger---where the
dense cropland proxy it replaced could never say a district had no observed fires.

\end{enumerate}



Together these bound the inventory as a transparent, reproducible \emph{backbone}
rather than a measured account: its magnitudes are auditable and its geography
reconciles by construction, but Tier-3 factors, a resolved rice water regime, and
district-level ground data --- each identified above --- are where further effort
would most improve it.
